\section{Circuit Complexity --- Tree NCAC}
%%% generated by the blueprint pipeline from TCSlib.Complexity.CircuitComplexity.TreeNCAC
\subsection*{Circuit Complexity --- Tree NCAC: Introduction}

Declarations of the Lean module \texttt{TCSlib.Complexity.CircuitComplexity.TreeNCAC}.

\subsection*{Circuit Complexity --- Tree NCAC: Results}

\begin{definition}[Family of Boolean tree circuits]
\lean{BoolCircuit.TreeCircuitFamily}
\label{BoolCircuit.TreeCircuitFamily}
\leanok
\uses{BoolCircuit.TreeCircuit}
A tree circuit family assigns to each input length $n$ a Boolean tree circuit on $n$
variables.
\end{definition}

\begin{definition}[Acceptance by a tree circuit family]
\lean{BoolCircuit.TreeCircuitFamily.Accepts}
\label{BoolCircuit.TreeCircuitFamily.Accepts}
\leanok
\uses{BoolCircuit.TreeCircuit.eval}
Let $C$ be a tree circuit family and let $w$ be a finite Boolean word. The family
accepts $w$ if its circuit for inputs of length $|w|$, evaluated on the entries of $w$,
outputs true.
\end{definition}

\begin{definition}[Language recognized by a circuit family]
\lean{BoolCircuit.TreeCircuitFamily.language}
\label{BoolCircuit.TreeCircuitFamily.language}
\leanok
\uses{BoolCircuit.TreeCircuitFamily.Accepts}
The language of a Boolean tree circuit family is the set of all Boolean words that the
family accepts.
\end{definition}

\begin{theorem}[Membership in a tree-circuit language]
\lean{BoolCircuit.TreeCircuitFamily.mem_language_iff}
\label{BoolCircuit.TreeCircuitFamily.mem_language_iff}
\leanok
\uses{BoolCircuit.TreeCircuit.eval, BoolCircuit.TreeCircuitFamily.language}
\difficulty{1}
For every Boolean tree-circuit family and every finite Boolean word $w$, the word $w$
belongs to the language recognized by the family if and only if the circuit assigned to
input length $|w|$ evaluates to true on the assignment given by the entries of $w$.
\end{theorem}

\begin{definition}[Polynomial size of a circuit family]
\lean{BoolCircuit.TreeCircuitFamily.IsPolySize}
\label{BoolCircuit.TreeCircuitFamily.IsPolySize}
\leanok
\uses{BoolCircuit.TreeCircuit.size}
A Boolean tree circuit family has polynomial size if there are natural numbers $a$ and
$k$ such that, for every input length $n$, the size of its circuit is at most
$a(n+1)^k$.
\end{definition}

\begin{definition}[Fan-in at most two]
\lean{BoolCircuit.TreeCircuitFamily.HasFaninTwo}
\label{BoolCircuit.TreeCircuitFamily.HasFaninTwo}
\leanok
\uses{BoolCircuit.TreeCircuit.maxFanin}
A Boolean tree circuit family has fan-in at most two if, for every input length $n$,
every gate in its circuit has at most two inputs.
\end{definition}

\begin{definition}[Polylogarithmic depth]
\lean{BoolCircuit.TreeCircuitFamily.HasPolylogDepth}
\label{BoolCircuit.TreeCircuitFamily.HasPolylogDepth}
\leanok
\uses{BoolCircuit.TreeCircuit.depth}
For a natural number $d$, a Boolean tree circuit family has depth bounded by a
degree-$d$ polynomial in $\log_2 n+1$ if there is a natural number $b$ such that, for
every input length $n$, its circuit depth is at most $b(\lfloor\log_2 n\rfloor+1)^d$.
\end{definition}

\begin{definition}[Tree circuit complexity class]
\lean{Language.InTreeNC}
\label{Language.InTreeNC}
\leanok
\uses{BoolCircuit.TreeCircuitFamily, BoolCircuit.TreeCircuitFamily.HasFaninTwo, BoolCircuit.TreeCircuitFamily.HasPolylogDepth, BoolCircuit.TreeCircuitFamily.IsPolySize, BoolCircuit.TreeCircuitFamily.language}
For a natural number $d$, a Boolean language belongs to the tree circuit class with
polynomial size and depth bounded by a degree-$d$ polynomial in $\lfloor\log_2 n\rfloor+1$ if
there is a family of Boolean tree circuits that recognizes the language, has fan-in at
most two at every input length, has polynomial size, and has depth at most a degree-$d$
polynomial in $\lfloor\log_2 n\rfloor+1$.
\end{definition}

\begin{definition}[Polynomial-size tree-AC language]
\lean{Language.InTreeAC}
\label{Language.InTreeAC}
\leanok
\uses{BoolCircuit.TreeCircuitFamily, BoolCircuit.TreeCircuitFamily.HasPolylogDepth, BoolCircuit.TreeCircuitFamily.IsPolySize, BoolCircuit.TreeCircuitFamily.language}
A language $L$ over the Boolean alphabet belongs to tree-AC at depth parameter $d$ if
there is a family of Boolean tree circuits of polynomial size, with polylogarithmic
depth parameterized by $d$, whose accepted language is exactly $L$.
\end{definition}

\begin{definition}[Tree-NC depth level]
\lean{BoolCircuit.TreeNCLevel}
\label{BoolCircuit.TreeNCLevel}
\leanok
\uses{Language.InTreeNC}
For a natural number $d$, the tree-NC level at parameter $d$ is the set of Boolean
languages that belong to tree-NC at parameter $d$.
\end{definition}

\begin{definition}[Tree-AC depth level]
\lean{BoolCircuit.TreeACLevel}
\label{BoolCircuit.TreeACLevel}
\leanok
\uses{Language.InTreeAC}
For a natural number $d$, the tree-AC level at parameter $d$ is the set of Boolean
languages that belong to tree-AC at parameter $d$.
\end{definition}

\begin{definition}[Tree-NC class]
\lean{BoolCircuit.TreeNC}
\label{BoolCircuit.TreeNC}
\leanok
\uses{BoolCircuit.TreeNCLevel}
The tree-NC class is the union, over all integer parameters $i \ge 1$, of the tree-NC
levels at parameter $i$.
\end{definition}

\begin{definition}[Tree-AC class]
\lean{BoolCircuit.TreeAC}
\label{BoolCircuit.TreeAC}
\leanok
\uses{BoolCircuit.TreeACLevel}
The tree-AC class is the union of the tree-AC levels over all natural-number parameters.
\end{definition}

\begin{theorem}[Characterization of tree NC]
\lean{BoolCircuit.mem_TreeNC_iff}
\label{BoolCircuit.mem_TreeNC_iff}
\leanok
\uses{BoolCircuit.TreeNC, BoolCircuit.TreeNCLevel, Language.InTreeNC}
\difficulty{1}
For every language $L$ over the Boolean alphabet, $L$ belongs to the union of the
tree-circuit classes with depth exponent at least one if and only if there is an integer
$i\ge 1$ and a family of Boolean circuit trees recognizing exactly $L$ whose fan-in is
at most two, whose size is polynomial in the input length, and whose depth is bounded by
a polynomial in $\log n+1$ of degree $i$.
\end{theorem}

\begin{theorem}[Characterization of tree AC]
\lean{BoolCircuit.mem_TreeAC_iff}
\label{BoolCircuit.mem_TreeAC_iff}
\leanok
\uses{BoolCircuit.TreeAC, BoolCircuit.TreeACLevel, Language.InTreeAC}
\difficulty{1}
For every language $L$ over the Boolean alphabet, $L$ belongs to the union of the
tree-circuit classes over all natural-number depth exponents if and only if there is a
natural number $i$ and a family of Boolean circuit trees recognizing exactly $L$ whose
size is polynomial in the input length and whose depth is bounded by a polynomial in
$\log n+1$ of degree $i$.
\end{theorem}

\begin{theorem}[Tree classes without fan-in bounds]
\lean{Language.InTreeNC.inTreeAC}
\label{Language.InTreeNC.inTreeAC}
\leanok
\uses{Language.InTreeAC, Language.InTreeNC}
\difficulty{4}
For a natural number $d$ and a language $L$ over Boolean strings, if $L$ is recognized
by a family of Boolean circuit trees with fan-in at most two, polynomial size, and depth
at most $b(\lfloor\log_2 n\rfloor+1)^d$ at every input length $n$, for some natural
number $b$, then $L$ is recognized by a family of Boolean circuit trees of polynomial
size and depth at most $b'(\lfloor\log_2 n\rfloor+1)^d$ at every input length $n$, for
some natural number $b'$.
\end{theorem}

\begin{definition}[Pair adjacent circuits]
\lean{BoolCircuit.pairUp}
\label{BoolCircuit.pairUp}
\leanok
\uses{BoolCircuit.TreeCircuit}
Given a Boolean $b$ and a list of Boolean tree circuits, form a new list by replacing
each consecutive pair of circuits with a single node whose gate is AND if $b$ is true
and OR if $b$ is false, and whose two children are that pair. Repeat this operation on
the remaining list; an empty list stays empty, and an unpaired final circuit is carried
over unchanged.
\end{definition}

\begin{theorem}[Length after pairing]
\lean{BoolCircuit.length_pairUp}
\label{BoolCircuit.length_pairUp}
\leanok
\uses{BoolCircuit.TreeCircuit, BoolCircuit.pairUp}
\difficulty{4}
For any Boolean gate type $b$ and any finite list $cs$ of Boolean circuit trees on $n$
variables, the list obtained by grouping consecutive circuits into pairs (leaving a
final unpaired circuit, if present) has length $\lceil |cs|/2\rceil$, equivalently
$(|cs|+1)/2$ with integer division.
\end{theorem}

\begin{theorem}[Evaluation after pairing children]
\lean{BoolCircuit.eval_node_pairUp}
\label{BoolCircuit.eval_node_pairUp}
\leanok
\uses{BoolCircuit.TreeCircuit, BoolCircuit.TreeCircuit.eval, BoolCircuit.pairUp}
\difficulty{4}
Let $n$ be a nonnegative integer, let $b$ be either true or false, let $x$ be an
assignment to the $n$ Boolean variables, and let $cs$ be a finite list of circuits.
Replacing consecutive pairs of children by gates of type $b$, leaving a final unpaired
child unchanged when the list has odd length, preserves the value of the gate of type
$b$ applied to the list.
\end{theorem}

\begin{theorem}[Depth bound for paired circuits]
\lean{BoolCircuit.maxDepth_pairUp}
\label{BoolCircuit.maxDepth_pairUp}
\leanok
\uses{BoolCircuit.TreeCircuit, BoolCircuit.TreeCircuit.depth, BoolCircuit.TreeCircuit.depth_node, BoolCircuit.TreeCircuit.maxDepth, BoolCircuit.TreeCircuit.maxDepth_cons, BoolCircuit.TreeCircuit.maxDepth_nil, BoolCircuit.pairUp}
\difficulty{4}
Let $n$ be a natural number, let $b$ specify either an AND or an OR gate, and let $cs$
be a list of circuits on $n$ variables. Pair consecutive circuits in $cs$ into gates of
type $b$, leaving a final unpaired circuit unchanged when the list has odd length. The
maximum depth in the resulting list is at most one plus the maximum depth in $cs$, where
the maximum over an empty list is $0$.
\end{theorem}

\begin{theorem}[Size bound under pairing]
\lean{BoolCircuit.sumSize_pairUp}
\label{BoolCircuit.sumSize_pairUp}
\leanok
\uses{BoolCircuit.TreeCircuit, BoolCircuit.TreeCircuit.size, BoolCircuit.TreeCircuit.size_node, BoolCircuit.TreeCircuit.sumSize, BoolCircuit.TreeCircuit.sumSize_cons, BoolCircuit.TreeCircuit.sumSize_nil, BoolCircuit.pairUp}
\difficulty{4}
Let $n$ be a natural number, let $b$ specify whether the new gates are AND or OR, and
let $cs$ be a finite list of Boolean circuits on $n$ variables. Pairing the circuits
according to the given pairing construction does not increase the sum of their sizes
plus the length of the list: the resulting sum of sizes plus resulting list length is at
most the original sum of sizes plus original list length.
\end{theorem}

\begin{theorem}[Fan-in bound for paired circuits]
\lean{BoolCircuit.maxFaninL_pairUp}
\label{BoolCircuit.maxFaninL_pairUp}
\leanok
\uses{BoolCircuit.TreeCircuit, BoolCircuit.TreeCircuit.maxFanin, BoolCircuit.TreeCircuit.maxFaninL, BoolCircuit.TreeCircuit.maxFaninL_cons, BoolCircuit.TreeCircuit.maxFaninL_nil, BoolCircuit.TreeCircuit.maxFanin_node, BoolCircuit.pairUp}
\difficulty{4}
Let $n$ be a natural number, let $b$ specify either an AND or an OR gate, and let $cs$
be a list of circuits on $n$ variables. Pair consecutive circuits in $cs$ into gates of
type $b$, leaving a final unpaired circuit unchanged when the list has odd length. The
maximum fan-in in the resulting list is at most the larger of $2$ and the maximum fan-in
in $cs$, where the maximum over an empty list is $0$.
\end{theorem}

\begin{definition}[Fuel-bounded circuit combination]
\lean{BoolCircuit.combineFuel}
\label{BoolCircuit.combineFuel}
\leanok
\uses{BoolCircuit.TreeCircuit, BoolCircuit.pairUp}
Given a Boolean $b$, a natural-number fuel value, and a list of Boolean tree circuits,
recursively combine circuits using the gate selected by $b$. With zero fuel, return one
node of that gate with the entire list as children; with positive fuel, return that gate
node on an empty list, return the sole circuit unchanged on a singleton list, and
otherwise pair adjacent circuits using that gate and recurse on the resulting list with
one less unit of fuel.
\end{definition}

\begin{definition}[Combine a list of circuits]
\lean{BoolCircuit.combine}
\label{BoolCircuit.combine}
\leanok
\uses{BoolCircuit.TreeCircuit, BoolCircuit.combineFuel}
Given a Boolean $b$ and a list of Boolean tree circuits, combine the list using the gate
selected by $b$, with fuel equal to the length of the list. The selected gate is AND
when $b$ is true and OR when $b$ is false.
\end{definition}

\begin{theorem}[Evaluation preserved by recursive combination]
\lean{BoolCircuit.combineFuel_eval}
\label{BoolCircuit.combineFuel_eval}
\leanok
\uses{BoolCircuit.TreeCircuit, BoolCircuit.TreeCircuit.eval, BoolCircuit.combineFuel, BoolCircuit.eval_node_pairUp, BoolCircuit.pairUp}
\difficulty{3}
Let $n$ be a natural number, let $b$ specify an AND or OR gate, let $x$ be an assignment
of Boolean values to the $n$ variables, and let $cs$ be a finite list of Boolean circuit
trees on those variables. For every natural number $k$, the circuit obtained by
recursively combining $cs$ with fuel $k$ has the same value under $x$ as the single gate
of type $b$ applied to the circuits in $cs$.
\end{theorem}

\begin{theorem}[Fanin bound for recursive circuit combination]
\lean{BoolCircuit.combineFuel_maxFanin}
\label{BoolCircuit.combineFuel_maxFanin}
\leanok
\uses{BoolCircuit.TreeCircuit, BoolCircuit.TreeCircuit.maxFanin, BoolCircuit.TreeCircuit.maxFaninL, BoolCircuit.TreeCircuit.maxFaninL_cons, BoolCircuit.TreeCircuit.maxFaninL_nil, BoolCircuit.TreeCircuit.maxFanin_node, BoolCircuit.combineFuel, BoolCircuit.length_pairUp, BoolCircuit.maxFaninL_pairUp, BoolCircuit.pairUp}
\difficulty{5}
Let $n$ be a natural number, let $b$ specify an AND or OR gate, and let $cs$ be a finite
list of Boolean circuit trees on $n$ variables. For every natural number $k$ with
$|cs|\le k$, the circuit obtained by recursively combining the list with fuel $k$ has
maximum fanin at most the larger of $2$ and the greatest fanin among circuits in $cs$,
taking that greatest fanin to be $0$ when the list is empty.
\end{theorem}

\begin{theorem}[Depth bound for recursive circuit combination]
\lean{BoolCircuit.combineFuel_depth}
\label{BoolCircuit.combineFuel_depth}
\leanok
\uses{BoolCircuit.TreeCircuit, BoolCircuit.TreeCircuit.depth, BoolCircuit.TreeCircuit.depth_node, BoolCircuit.TreeCircuit.maxDepth, BoolCircuit.TreeCircuit.maxDepth_cons, BoolCircuit.TreeCircuit.maxDepth_nil, BoolCircuit.combineFuel, BoolCircuit.length_pairUp, BoolCircuit.maxDepth_pairUp, BoolCircuit.pairUp}
\difficulty{6}
Let $n$ be a natural number, let $b$ specify an AND or OR gate, and let $cs$ be a finite
list of Boolean circuit trees on $n$ variables. For every natural number $k$ with
$|cs|\le k$, the circuit obtained by recursively combining the list with fuel $k$ has
depth at most the maximum depth in $cs$ plus $\operatorname{clog}_2(|cs|)+1$, where the
maximum over an empty list is $0$.
\end{theorem}

\begin{theorem}[Size bound for recursive circuit combination]
\lean{BoolCircuit.combineFuel_size}
\label{BoolCircuit.combineFuel_size}
\leanok
\uses{BoolCircuit.TreeCircuit, BoolCircuit.TreeCircuit.size, BoolCircuit.TreeCircuit.size_node, BoolCircuit.TreeCircuit.sumSize, BoolCircuit.TreeCircuit.sumSize_cons, BoolCircuit.TreeCircuit.sumSize_nil, BoolCircuit.combineFuel, BoolCircuit.pairUp, BoolCircuit.sumSize_pairUp}
\difficulty{4}
Let $n$ be a natural number, let $b$ specify an AND or OR gate, and let $cs$ be a finite
list of Boolean circuit trees on $n$ variables. For every natural number $k$, the
circuit obtained by recursively combining $cs$ with fuel $k$ has size at most the sum of
the sizes of the circuits in $cs$, plus $|cs|+1$.
\end{theorem}

\begin{theorem}[Evaluation of a combined circuit]
\lean{BoolCircuit.combine_eval}
\label{BoolCircuit.combine_eval}
\leanok
\uses{BoolCircuit.TreeCircuit, BoolCircuit.TreeCircuit.eval, BoolCircuit.combine, BoolCircuit.combineFuel_eval}
\difficulty{1}
For any number $n$ of variables, any Boolean gate type, any finite list $cs$ of Boolean
circuit trees on those variables, and any assignment $x$ of Boolean values to the
variables, combining $cs$ gives the same value as applying a gate of the chosen type
directly to the list of circuits in $cs$.
\end{theorem}

\begin{theorem}[Fanin of a combined circuit]
\lean{BoolCircuit.combine_maxFanin}
\label{BoolCircuit.combine_maxFanin}
\leanok
\uses{BoolCircuit.TreeCircuit, BoolCircuit.TreeCircuit.maxFanin, BoolCircuit.TreeCircuit.maxFaninL, BoolCircuit.combine, BoolCircuit.combineFuel_maxFanin}
\difficulty{1}
For any number $n$ of variables, any Boolean gate type, and any finite list $cs$ of
Boolean circuit trees on those variables, the maximum fanin of the circuit obtained by
combining the list is at most the maximum of $2$ and the largest fanin among the
circuits in $cs$, where the largest fanin of an empty list is $0$.
\end{theorem}

\begin{theorem}[Depth of a combined circuit]
\lean{BoolCircuit.combine_depth}
\label{BoolCircuit.combine_depth}
\leanok
\uses{BoolCircuit.TreeCircuit, BoolCircuit.TreeCircuit.depth, BoolCircuit.TreeCircuit.maxDepth, BoolCircuit.combine, BoolCircuit.combineFuel_depth}
\difficulty{1}
For any number $n$ of variables, any Boolean gate type, and any finite list $cs$ of
Boolean circuit trees on those variables, the depth of the circuit obtained by combining
the list is at most the maximum depth in $cs$, plus $\lceil\log_2 |cs|\rceil+1$, where
the maximum depth of an empty list is $0$.
\end{theorem}

\begin{theorem}[Size of a combined circuit]
\lean{BoolCircuit.combine_size}
\label{BoolCircuit.combine_size}
\leanok
\uses{BoolCircuit.TreeCircuit, BoolCircuit.TreeCircuit.size, BoolCircuit.TreeCircuit.sumSize, BoolCircuit.combine, BoolCircuit.combineFuel_size}
\difficulty{1}
For any number $n$ of variables, any Boolean gate type, and any finite list $cs$ of
Boolean circuit trees on those variables, the size of the circuit obtained by combining
the list is at most the sum of the sizes of the circuits in $cs$, plus the length of
$cs$, plus $1$.
\end{theorem}

\begin{definition}[Binary-arity circuit conversion]
\lean{BoolCircuit.TreeCircuit.toBinary}
\label{BoolCircuit.TreeCircuit.toBinary}
\leanok
\uses{BoolCircuit.TreeCircuit, BoolCircuit.combine}
For every Boolean circuit tree on $n$ variables, define its binary conversion
recursively: literals remain unchanged, while each gate keeps its AND or OR type,
converts every child, and combines the converted children in successive pairwise layers
until at most one circuit remains.
\end{definition}

\begin{theorem}[Evaluation under gate-preserving replacement]
\lean{BoolCircuit.eval_node_map}
\label{BoolCircuit.eval_node_map}
\leanok
\uses{BoolCircuit.TreeCircuit, BoolCircuit.TreeCircuit.eval}
\difficulty{5}
Let $n$ be a nonnegative integer, let $b$ be either true or false, let $x$ be an
assignment to the $n$ Boolean variables, and let $f$ map circuits to circuits. If $f(c)$
and $c$ have the same value under $x$ for every circuit $c$ in a finite list, then
replacing each list entry $c$ by $f(c)$ leaves the value of the gate with connective $b$
unchanged.
\end{theorem}

\begin{theorem}[Depth bound for mapped circuits]
\lean{BoolCircuit.maxDepth_map_le}
\label{BoolCircuit.maxDepth_map_le}
\leanok
\uses{BoolCircuit.TreeCircuit, BoolCircuit.TreeCircuit.depth, BoolCircuit.TreeCircuit.maxDepth, BoolCircuit.TreeCircuit.maxDepth_cons}
\difficulty{4}
Let $n$ and $m$ be natural numbers, let $f$ map circuits on $n$ variables to circuits on
$n$ variables, and let $cs$ be a list of such circuits. If every $c\in cs$ satisfies
$\mathrm{depth}(f(c))\le m$, then the maximum depth among the circuits $f(c)$ for $c\in
cs$ is at most $m$, with the maximum defined to be $0$ for an empty list.
\end{theorem}

\begin{theorem}[Fan-in bound under mapping]
\lean{BoolCircuit.maxFaninL_map_le}
\label{BoolCircuit.maxFaninL_map_le}
\leanok
\uses{BoolCircuit.TreeCircuit, BoolCircuit.TreeCircuit.maxFanin, BoolCircuit.TreeCircuit.maxFaninL, BoolCircuit.TreeCircuit.maxFaninL_cons}
\difficulty{4}
Let $n,m\in\mathbb{N}$, let $f$ map Boolean circuit trees on $n$ variables to Boolean
circuit trees on $n$ variables, and let $cs$ be a finite list of such trees. Suppose
every tree $c$ occurring in $cs$ has maximum fan-in at most $m$ after applying $f$. Then
the maximum, over the images $f(c)$ for $c$ in the list, of their recursively defined
fan-in bounds is at most $m$, where the maximum for an empty list is $0$.
\end{theorem}

\begin{theorem}[Evaluation preserved by binary conversion]
\lean{BoolCircuit.toBinary_eval}
\label{BoolCircuit.toBinary_eval}
\leanok
\uses{BoolCircuit.TreeCircuit, BoolCircuit.TreeCircuit.eval, BoolCircuit.TreeCircuit.ind, BoolCircuit.TreeCircuit.toBinary, BoolCircuit.combine_eval, BoolCircuit.eval_node_map}
\difficulty{4}
For every natural number $n$, every Boolean circuit tree $c$ on $n$ variables, and every
assignment $x : \bbr^n \to \{0,1\}$, the binary conversion of $c$ has the same evaluation
at $x$ as $c$.
\end{theorem}

\begin{theorem}[Fan-in after conversion to binary]
\lean{BoolCircuit.toBinary_maxFanin_le}
\label{BoolCircuit.toBinary_maxFanin_le}
\leanok
\uses{BoolCircuit.TreeCircuit, BoolCircuit.TreeCircuit.ind, BoolCircuit.TreeCircuit.maxFanin, BoolCircuit.TreeCircuit.toBinary, BoolCircuit.combine_maxFanin, BoolCircuit.maxFaninL_map_le}
\difficulty{4}
For every Boolean circuit tree $c$, the maximum fan-in of its binary conversion is at
most $2$.
\end{theorem}

\begin{theorem}[Aggregate size bound for binary conversion]
\lean{BoolCircuit.sumSize_map_toBinary}
\label{BoolCircuit.sumSize_map_toBinary}
\leanok
\uses{BoolCircuit.TreeCircuit, BoolCircuit.TreeCircuit.size, BoolCircuit.TreeCircuit.sumSize, BoolCircuit.TreeCircuit.sumSize_cons, BoolCircuit.TreeCircuit.sumSize_nil, BoolCircuit.TreeCircuit.toBinary}
\difficulty{4}
Let $n$ be a natural number and let $cs$ be a list of Boolean circuit trees on $n$
variables. Suppose that for every circuit $c$ in $cs$, the size of its binary conversion
plus one is at most three times its original size. Then the sum of the sizes of the
binary conversions of the circuits in $cs$, plus the length of $cs$, is at most three
times the sum of their original sizes.
\end{theorem}

\begin{theorem}[Strict size bound after binary conversion]
\lean{BoolCircuit.toBinary_size_succ_le}
\label{BoolCircuit.toBinary_size_succ_le}
\leanok
\uses{BoolCircuit.TreeCircuit, BoolCircuit.TreeCircuit.ind, BoolCircuit.TreeCircuit.size, BoolCircuit.TreeCircuit.size_node, BoolCircuit.TreeCircuit.toBinary, BoolCircuit.combine_size, BoolCircuit.sumSize_map_toBinary}
\difficulty{4}
For every Boolean circuit tree $c$, the size of its binary conversion plus one is at
most three times the size of $c$.
\end{theorem}

\begin{theorem}[Size after conversion to binary]
\lean{BoolCircuit.toBinary_size_le}
\label{BoolCircuit.toBinary_size_le}
\leanok
\uses{BoolCircuit.TreeCircuit, BoolCircuit.TreeCircuit.size, BoolCircuit.TreeCircuit.toBinary, BoolCircuit.toBinary_size_succ_le}
\difficulty{1}
For every Boolean circuit tree $c$, the size of its binary conversion is at most three
times the size of $c$.
\end{theorem}

\begin{theorem}[Depth after conversion to binary]
\lean{BoolCircuit.toBinary_depth_le}
\label{BoolCircuit.toBinary_depth_le}
\leanok
\uses{BoolCircuit.TreeCircuit, BoolCircuit.TreeCircuit.depth, BoolCircuit.TreeCircuit.depth_le_maxDepth, BoolCircuit.TreeCircuit.depth_node, BoolCircuit.TreeCircuit.ind, BoolCircuit.TreeCircuit.maxDepth, BoolCircuit.TreeCircuit.maxFanin, BoolCircuit.TreeCircuit.maxFaninL, BoolCircuit.TreeCircuit.maxFanin_le_maxFaninL, BoolCircuit.TreeCircuit.maxFanin_node, BoolCircuit.TreeCircuit.toBinary, BoolCircuit.combine_depth, BoolCircuit.maxDepth_map_le}
\difficulty{5}
Let $c$ be a Boolean circuit tree and let $w$ be a nonnegative integer. If every gate of
$c$ has fan-in at most $w$, then the depth of its binary conversion is at most
$\mathrm{depth}(c)\bigl(\lceil\log_2 w\rceil+1\bigr)$, where the ceiling logarithm is
the least integer $k$ such that $2^k\ge w$.
\end{theorem}

\begin{theorem}[Logarithmic bound for a polynomial expression]
\lean{BoolCircuit.clog_poly_le}
\label{BoolCircuit.clog_poly_le}
\leanok
\difficulty{3}
For all natural numbers $a$, $k$, and $m$, the base-$2$ ceiling logarithm satisfies
$\operatorname{clog}_2\!\bigl(a(m+1)^k\bigr)\le
a+k\bigl(\operatorname{log}_2(m)+1\bigr)$.
\end{theorem}

\begin{theorem}[Tree classes and fan-in-two depth]
\lean{Language.InTreeAC.inTreeNC_succ}
\label{Language.InTreeAC.inTreeNC_succ}
\leanok
\uses{BoolCircuit.TreeCircuit.depth, BoolCircuit.TreeCircuit.maxFanin, BoolCircuit.TreeCircuit.maxFanin_le_size, BoolCircuit.TreeCircuit.size, BoolCircuit.TreeCircuit.toBinary, BoolCircuit.TreeCircuitFamily.mem_language_iff, BoolCircuit.clog_poly_le, BoolCircuit.toBinary_depth_le, BoolCircuit.toBinary_eval, BoolCircuit.toBinary_maxFanin_le, BoolCircuit.toBinary_size_le, Language.InTreeAC, Language.InTreeNC}
\difficulty{5}
For a natural number $d$ and a language $L$ over Boolean strings, if $L$ is recognized
by a family of Boolean circuit trees of polynomial size and depth at most
$b(\lfloor\log_2 n\rfloor+1)^d$ at every input length $n$, for some natural number $b$,
then $L$ is recognized by such a tree family with fan-in at most two and polynomial
size, and with depth at most $b'(\lfloor\log_2 n\rfloor+1)^{d+1}$ at every input length
$n$, for some natural number $b'$.
\end{theorem}

\begin{theorem}[Tree circuit complexity class inclusion]
\lean{BoolCircuit.TreeNCLevel_subset_ACLevel}
\label{BoolCircuit.TreeNCLevel_subset_ACLevel}
\leanok
\uses{BoolCircuit.TreeACLevel, BoolCircuit.TreeNCLevel, Language.InTreeNC.inTreeAC}
\difficulty{1}
For every natural number $i$, every language recognized by a polynomial-size
tree-circuit family of polylogarithmic depth exponent $i$ and fan-in at most two belongs
to the class of languages recognized by polynomial-size tree-circuit families of
polylogarithmic depth exponent $i$.
\end{theorem}

\begin{theorem}[Tree circuit complexity class inclusion]
\lean{BoolCircuit.TreeACLevel_subset_NCLevel_succ}
\label{BoolCircuit.TreeACLevel_subset_NCLevel_succ}
\leanok
\uses{BoolCircuit.TreeACLevel, BoolCircuit.TreeNCLevel, Language.InTreeAC.inTreeNC_succ}
\difficulty{1}
For every natural number $i$, every language recognized by a family of Boolean tree
circuits whose size is bounded by a polynomial and whose depth is at most $b(1+\log_2
n)^i$ for some constant $b$ at every input length $n$ is also recognized by such a
family with fan-in at most two and depth at most $b'(1+\log_2 n)^{i+1}$ for some
constant $b'$.
\end{theorem}

\begin{theorem}[Equality of tree circuit classes]
\lean{BoolCircuit.TreeNC_eq_TreeAC}
\label{BoolCircuit.TreeNC_eq_TreeAC}
\leanok
\uses{BoolCircuit.TreeAC, BoolCircuit.TreeNC, BoolCircuit.mem_TreeAC_iff, BoolCircuit.mem_TreeNC_iff, Language.InTreeAC.inTreeNC_succ, Language.InTreeNC.inTreeAC}
\difficulty{4}
The class formed by taking the union, over all integers $i\ge 1$, of languages
recognized by polynomial-size tree-circuit families with fan-in at most two and depth at
most $b(\log_2 n+1)^i$ for some constant $b$ is equal to the class formed by taking the
union over all natural numbers $i$ of languages recognized by polynomial-size
tree-circuit families with depth at most $b(\log_2 n+1)^i$ for some constant $b$.
\end{theorem}
